% from ega4-11-fr.tex
\oldpage[IV-3]{131}

d'où, par comparaison avec (11.2.9.7), un isomorphisme de $C'$-modules
gradués
\[
Q\otimes_AA'\xrightarrow{\sim}Q_0\otimes_{A'_0}A'.
\]
Comme $C'_0$ est noethérien, $Q_0$ est un $C'_0$-module de type fini, donc
$Q_0\otimes_{A'_0}A'$ est un $C'$-module gradué de type fini, et il en est par
suite de même de $Q\otimes_AA'$~; soient $q_h$ ($1\leq h\leq s$) des éléments
homogènes de $Q$ dont les images dans $Q\otimes_AA'$ engendrent ce
$C'$-module.

Considérons maintenant un idéal premier quelconque
$\mathfrak p\supset\mathfrak RB$ de $B$~; on a
$C_\mathfrak p=\operatorname{gr}_{\mathfrak J_\mathfrak p}^*(B_\mathfrak p)$
par platitude, et
$\operatorname{gr}_{\mathfrak J_\mathfrak p}^0(B_\mathfrak p)
=B_\mathfrak p/\mathfrak J_\mathfrak p$ est réduit à 0 ou est un anneau
local~; pour prouver le lemme, on peut évidemment se borner au second cas. Il
est clair alors que chacun des $(B_\mathfrak p/\mathfrak J_\mathfrak p)$-modules
$\operatorname{gr}_{\mathfrak J_\mathfrak p}^i(B_\mathfrak p)$ est de
présentation finie. D'autre part, pour tout indice $i$,
$M/\mathfrak J^iM$ s'identifie à
$(M_0/\mathfrak J_0^iM_0)\otimes_{A_0}A$, et est donc un $B$-module de
présentation finie. Par récurrence sur $i$, l'hypothèse que
$\operatorname{gr}_{\mathfrak J}^*(M)$ est un $A$-module plat entraîne, en
vertu de ($\mathbf 0_{\mathrm I}$, 6.1.2), que $M/\mathfrak J^iM$ est un
$A$-module plat. L'application de (11.2.6.1) où l'on remplace $M_0$ par
$M_0/\mathfrak J_0^kM_0$ pour $k\leq i$ montre qu'il existe un indice
$\lambda_i$ tel que, pour $\mu\geq\lambda_i$, chacun des
$M_\mu/\mathfrak J_\mu^{k+1}M_\mu$ soit un $A_\mu$-module plat, et par suite
aussi ($\mathbf 0_{\mathrm I}$, 6.1.2) chacun des
$\operatorname{gr}_{\mathfrak J_\mu}^k(M_\mu)$ est un $A_\mu$-module plat
pour $k\leq i$. On déduit par suite de (11.2.9.2) que chacun des
$(B/\mathfrak J)$-modules $\operatorname{gr}_{\mathfrak J}^i(M)$ est de
présentation finie, donc
$\operatorname{gr}_{\mathfrak J_\mathfrak p}^i(M_\mathfrak p)$ est un
$(B_\mathfrak p/\mathfrak J_\mathfrak p)$-module de présentation finie. Par
ailleurs les images des $q_h$ dans
\[
Q_\mathfrak p\otimes_{B_\mathfrak p}k(\mathfrak p)
=(Q\otimes_AA')_\mathfrak p\otimes_{A'}k(\mathfrak p)
\]
engendrent ce module sur
$C_\mathfrak p\otimes_{B_\mathfrak p}k(\mathfrak p)$, puisque
\[
C_\mathfrak p\otimes_{B_\mathfrak p}k(\mathfrak p)
=C'_\mathfrak p\otimes_{A'}k(\mathfrak p)~;
\]
comme on a la suite exacte
\[
Q_\mathfrak p\otimes_{B_\mathfrak p}k(\mathfrak p)
\longrightarrow C_\mathfrak p^{r}\otimes_{B_\mathfrak p}k(\mathfrak p)
\longrightarrow
\operatorname{gr}_{\mathfrak J_\mathfrak p}^*(M_\mathfrak p)
\otimes_{B_\mathfrak p}k(\mathfrak p)\longrightarrow0,
\]
on en conclut que
$\operatorname{gr}_{\mathfrak J_\mathfrak p}^*(M_\mathfrak p)
\otimes_{B_\mathfrak p}k(\mathfrak p)$ est un
$(C_\mathfrak p\otimes_{B_\mathfrak p}k(\mathfrak p))$-module de présentation
finie. Appliquant le lemme (11.2.9.4) où $A$, $B$, $M$ sont remplacés
respectivement par $B_\mathfrak p$, $C_\mathfrak p$ et
$\operatorname{gr}_{\mathfrak J_\mathfrak p}^*(M_\mathfrak p)$, on en conclut
que $Q_\mathfrak p$ est un $C_\mathfrak p$-module de type fini. Utilisant
maintenant le lemme de Nakayama (et le fait que $C=B$) on en déduit que les
images des $q_h$ dans $Q_\mathfrak p$ engendrent ce $C_\mathfrak p$-module.
Ceci achève de prouver le lemme (11.2.9.6) et le fait que
$\operatorname{gr}_{\mathfrak J}^*(M)$ est un $C$-module de présentation
finie.

\medskip
\noindent\textbf{VI)} \emph{Fin de la démonstration.} --- Posons pour abréger
$C_\lambda=C_0\otimes_{A_0}A_\lambda$ (égal en fait à $B_\lambda$),
$N_\lambda=\operatorname{gr}_{\mathfrak J_\lambda}^*(M_\lambda)$ et
$N=\operatorname{gr}_{\mathfrak J}^*(M)$. Notons d'abord que pour chaque
entier $k$, $\mathfrak J^kM$ s'identifie à
$\varinjlim\mathfrak J_\lambda^kM_\lambda$ par exactitude du foncteur
$\varinjlim$~; l'image de $\mathfrak J_\lambda^kM_\lambda$ dans $M_\mu$ pour
$\lambda\leq\mu$ (resp. dans $M$) engendrant $\mathfrak J_\mu^kM_\mu$
(resp. $\mathfrak J^kM$)~; utilisant encore l'exactitude de $\varinjlim$, on
en conclut que $N$ s'identifie canoniquement à $\varinjlim N_\lambda$. Faisant
cette identification, nous allons d'abord prouver que~:
\begin{equation}
\text{\emph{Pour $\lambda$ assez grand, l'homomorphisme canonique
$v_\lambda:N_\lambda\otimes_{A_\lambda}A\to N$ est bijectif.}}
\label{IV.11.2.9.9-fr}
\tag{11.2.9.9}
\end{equation}
Comme les $C_\lambda$ sont noethériens et $N_\lambda$ un $C_\lambda$-module
de type fini, les $C$-modules $N_\lambda\otimes_{A_\lambda}A$ sont de
présentation finie et forment un système inductif filtrant, dont la limite
inductive s'identifie canoniquement à $N$ en vertu du fait que $\varinjlim$
commute aux produits tensoriels. En outre, les homomorphismes de transition
$v_{\mu\lambda}:N_\lambda\otimes_{A_\lambda}A
\to N_\mu\otimes_{A_\mu}A$

